\section{Scientific notation}

Physics and science in general often deal with very large and very small numbers. Writing such numbers in ordinary decimal notation is inconvenient and often hides the scale of the quantity. Scientific notation makes the scale visible.

\begin{definitionbox}
A number is written in scientific notation if it has the form
\[
a \times 10^n,
\]
where \(n\) is an integer and usually \(1 \leq |a| < 10\).
\end{definitionbox}

The number \(a\) carries the significant digits, while \(10^n\) carries the scale. For example,
\[
2.7\times 10^{17}=270000000000000000,
\]
whereas
\[
1.602\times 10^{-19}=0.0000000000000000001602.
\]
The sign of the exponent tells us whether the decimal point is moved to the right or to the left.

This notation is not optional. It is a basic part of scientific language. Without it, many numbers used in science would be almost unreadable.

The elementary charge is approximately
\[
e = 1.602 \times 10^{-19}\,\mathrm{C},
\]
and the Avogadro constant is approximately
\[
N_A = 6.022 \times 10^{23}\,\mathrm{mol}^{-1}.
\]
The notation immediately shows that these quantities belong to completely different scales.

Scientific notation also makes multiplication and division easier. We use the rules
\[
10^a \cdot 10^b = 10^{a+b},
\qquad
\frac{10^a}{10^b} = 10^{a-b}.
\]

\begin{examplebox}
\[
(2 \times 10^3)(3 \times 10^4)
= 6 \times 10^{3+4}
= 6 \times 10^7.
\]
Also,
\[
\frac{8 \times 10^6}{2 \times 10^2}
= 4 \times 10^{6-2}
= 4 \times 10^4.
\]
\end{examplebox}

\begin{remarkbox}
The exponent of ten gives the order of magnitude. In many scientific situations, understanding the order of magnitude is more important than knowing many decimal digits.
\end{remarkbox}

The reader is expected to understand notation such as
\[
10^{-6}, \qquad 10^3, \qquad 2.7 \times 10^{17}, \qquad 6.022 \times 10^{23}.
\]
